\paragraph{Borderline items.} We deliberately filter each dataset to its difficult, borderline items
(Appendix~\ref{sec:appendix-sampling}). On an unfiltered, naturally distributed 100-item WildGuard binary sample,
L1--L5 wiggle rates are 5.3--12.7 percentage points lower than on the selected hard sample, while L6 coverage is nearly
unchanged (70.3\% versus 69.7\%; Appendix~\ref{sec:appendix-wildguard-representative}). This confirms that hard-item
selection inflates absolute rates under low-to-moderate pressure, while providing initial evidence that the L6 result is
not solely an artifact of that selection. The ablation covers only one dataset and one grading scale; other datasets and
scales still require equivalent controls.

\paragraph{No human baseline.} We do not measure human-annotator wiggle under the same settings,
and thus we cannot establish the relationship between the wiggle of LLMs and humans. WildGuard's
human-consensus data suggests that there may be some 
non-trivial connection between judge wiggle and inter-human disagreement (Appendix~\ref{sec:appendix-human-consensus}), but a full baseline human study
under the L1--L6 ladder would be needed to settle this. Existing work shows that
LLM-generated arguments can shift human opinions and that LLMs can be persuasive in multi-turn debates with
humans~\citep{durmus2024persuasiveness,salvi2025gpt4}, but it is unclear how these findings transfer to the judging
domains and pressure protocols of our study.

\paragraph{Dataset sample sizes.} We sample 100 items from each dataset per grading scale (50 prompt pairs for Paired Prompts). 
The pooled per-level and per-judge results are stable enough to support our qualitative conclusions, but 
per-cell point estimates for more granular slices of data like in the appendices should be interpreted with appropriate caution.

\paragraph{Single L6 persuader set.} The L6 adaptive persuader pool is fixed at three models (GPT-5.4, Claude 4.6 Opus, 
Grok-4.1 Reasoning). The models come from different organizations with capabilities assumed sufficient for adaptive argumentation. 
Different persuader configurations like using a larger set of models or including models fine-tuned for adversarial persuasion could 
produce different wiggle and self-persuasion patterns (Appendix~\ref{sec:appendix-self}).

\paragraph{Dataset coverage.} The six datasets span safety, toxicity, AI-text detection, and political-content evaluation, 
but they do not exhaust the space of LLM-as-judge applications. Aesthetic judgment, code-review correctness, 
mathematical-reasoning verification, medical-content review, and other expert-evaluation tasks may produce 
qualitatively different wiggle profiles. Why some tasks are more or less robust than others is an open question for future work.

\paragraph{Causal claims on flip asymmetries.} Our results are correlational. We observe that pressure changes verdicts and that the 
\emph{direction} of those changes may be consistent with what one would expect from training priors 
(\S\ref{fo:structure} with binary flips leaning restrictive, Likert flips leaning permissive), but we cannot definitively 
establish the causal mechanism. Alternative explanations like inherent task asymmetries, rubric-prompt 
construction effects, or systematic differences in argument quality between the two flip directions could factor 
into the same observations and cannot be fully ruled out without additional controlled experiments.

\paragraph{Black-box techniques.} The Wiggle Framework deliberately emphasizes black-box methods, requiring only observed verdicts. 
White-box or other mechanistic approaches that examine attention, residual streams, logits, or other activation patterns might uncover
additional signal for predicting when certain models or items wiggle.